\subsection{Strict Output Enforcement Protocols}

In order to guarantee the machine readability and stability of the model's outputs, the framework delineates three standardised interaction protocols, with the objective of accommodating varying model capabilities and application scenarios.

\begin{enumerate}
\item \textbf{Direct Mode}: The model is compelled to produce solely the words \texttt{true} or \texttt{false}. This mode has been engineered to maximally suppress autoregressive interference from Chain of Thought (CoT) reasoning, rendering it suitable for large-scale validation scenarios requiring extremely high throughput and low latency.
\item \textbf{SSE Mode}: The data is arranged in a row-based \texttt{key: value} format, for example, \texttt{reason: ... \textbackslash n answer: ...}. This approach maintains the interpretability of the reasoning process while supporting streaming parsing, facilitating real-time monitoring and evaluation of the logic.
\item \textbf{JSON Mode}: The implementation of complete JSON object output is enforced, rendering it suitable for scenarios necessitating programmatic parsing of complex reasoning logic or integration with other automated systems.
\end{enumerate}

The aforementioned three modes collectively constitute the framework's interface layer. By means of a strict constraint on the output space through prompt engineering, the transformation of natural language-generated evaluation tasks into quasi-deterministic classification tasks is achieved.
